\section{Methods}
\label{sec:methods}

\subsection{Notation and Setup}

We work on the 2D torus GridWorld from~\citet{rambaud2025mapformer}. An
environment of size $N \times N$ (typically $N{=}64$) assigns to each cell
an observation type from $K + 1$ symbols: $K$ regular types plus a blank
token $B$. A fraction $p_\mathrm{empty}$ of cells are blank. Optionally,
$L$ additional cells carry unique single-use landmark tokens (disjoint
from the $K$ regular types).

The agent emits an interleaved trajectory
\begin{equation}
  s = (a_1, o_1, a_2, o_2, \ldots, a_T, o_T)
\end{equation}
with actions $a_t \in \{\uparrow,\downarrow,\leftarrow,\rightarrow\}$
sampled over directed walks ($k$ steps per chosen direction,
$k \in \{1,\ldots,10\}$) and observations $o_t$ determined by the
obs-map at the agent's current cell after applying $a_t$. We train on
$\sim$1M sequences ($T{=}128$ steps, batch size 128, AdamW
$\mathrm{lr}{=}3\!\times\!10^{-4}$, weight decay 0.05, linear LR decay).

\paragraph{Task.} Self-supervised next-token prediction with loss applied
only at \emph{revisited} observation positions, matching MapFormer's
forced-navigation task: ``predict the upcoming observation each time the
agent returns to a previously visited location.''

\subsection{MapFormer Baseline (reproduction)}

A 1-layer Transformer on the unified vocabulary (actions + observations).
For each token $x_t$, a low-rank projection produces the Lie-algebra
increment $\Delta_t \in \mathbb{R}^{H \times B}$, where $H$ is the number
of attention heads and $B$ the number of $2{\times}2$ rotation blocks per
head. Angular velocities $\omega \in \mathbb{R}^{H \times B}$ are
learnable with geometric initialization
\begin{equation}
  \omega_i = \omega_{\max} \cdot \left( \tfrac{1}{\Delta_{\max}} \right)^{(i-1)/(B-1)}, \quad i = 1, \ldots, B
\end{equation}
(The paper's Eq.~17 specifies
$\omega_i = \omega_{\max} \cdot (1/\Delta_{\max})^{-i/B}$; the negated
exponent is a sign typo, confirmed by the paper's own inequality
$\omega_\mathrm{min} = \omega_{\max}/\Delta_{\max}$ and RoPE's analogous
decreasing schedule.)

Path integration proceeds in the Lie algebra via parallel cumulative sum:
\begin{equation}
  \theta_t = \omega \odot \mathrm{cumsum}_t(\Delta)
\end{equation}
where $\odot$ is elementwise product. The cumulative angles are applied to
queries and keys via RoPE. Attention is standard causal self-attention.

Correct reproduction requires seven specific fixes (torus boundary,
unified action/observation vocabulary, revisit-only loss, $\omega$ sign,
Hadamard EM attention, separate $q_0^p/k_0^p$, cumsum+exp not
prefix-product). With all fixes, our reproduction matches the paper's
reported 95.5\% accuracy for MapFormer-WM at $T{=}128$.

\subsection{Level 1 InEKF: Parallel Invariant EKF on SO(2)}

The Level~1 variant adds explicit state correction via a Kalman filter.
Following~\citet{markovic2017wrapping}, on $\mathrm{SO}(2)$ the
wrapped-innovation EKF is mathematically identical to the Lie-Group EKF,
so we operate on an unbounded $\hat\theta \in \mathbb{R}$ with
innovations wrapped to $[-\pi, \pi]$ via $\mathrm{atan2}(\sin, \cos)$.

A learned measurement head produces $\hat z_t \in [-\pi, \pi]$ from the
content embedding:
\begin{equation}
  \hat z_t = \pi \cdot \tanh\bigl( W_2 \cdot \mathrm{GELU}(W_1 \cdot x_t) \bigr)
\end{equation}

With constant process noise $Q$ and constant measurement noise $R$, the
scalar discrete algebraic Riccati equation has the closed-form fixed
point:
\begin{equation}
  P^* = \frac{-Q + \sqrt{Q^2 + 4QR}}{2}, \qquad K^* = \frac{P^* + Q}{P^* + Q + R}
\end{equation}

With $K^*$ frozen, the correction recurrence becomes a scalar affine
recurrence with constant coefficient:
\begin{equation}
  d_t = (1 - K^*) \, d_{t-1} + K^* \cdot \nu_t,
  \qquad
  \nu_t = \mathrm{atan2}\bigl( \sin(\hat z_t - \theta_{\mathrm{path},t}), \cos(\hat z_t - \theta_{\mathrm{path},t}) \bigr)
\end{equation}
which is equivalent to a causal 1-D convolution with kernel
$K^* (1 - K^*)^k,\, k = 0, 1, 2, \ldots$ computable via FFT in
$\mathcal{O}(T \log T)$ work / $\mathcal{O}(\log T)$ depth. The corrected
angles $\hat\theta_t = \theta_{\mathrm{path},t} + d_t$ are fed to RoPE in
place of $\theta_\mathrm{path}$.

\subsection{Level 2 InEKF: Full Heteroscedastic}

Level~2 learns a per-token measurement noise $R_t$ from the content
embedding and tracks a time-varying covariance $\Pi_t$. The covariance
recurrence under predict + update is
\begin{equation}
  \Pi_{t+1} = \frac{(\Pi_t + Q) \cdot R_t}{\Pi_t + Q + R_t}
\end{equation}
which is a M\"obius transformation $\Pi_{t+1} = M_t(\Pi_t)$ with matrix
\begin{equation}
  M_t = \begin{bmatrix} R_t + Q & Q \cdot R_t \\ 1 & R_t \end{bmatrix}
\end{equation}
M\"obius transforms compose under matrix multiplication, so a parallel
associative scan of $2 \times 2$ matrices yields $\Pi_t$ for all $t$ in
$\mathcal{O}(\log T)$ depth. The gain is then $K_t = \Pi_t / R_t$ (a
scalar simplification derivable from the predict--update chain). The
state correction uses a second affine scan with time-varying
$\alpha_t = 1 - K_t$.

Two parallel $2 \times 2$ matrix scans are $\sim 60\times$ slower in
wall-clock than Level~1's FFT convolution at our sequence lengths,
despite equivalent asymptotics. More importantly, Level~2 \emph{performs
worse} than Level~1.5 despite being strictly more expressive
(Section~\ref{sec:results}).

\subsection{Level 1.5 InEKF: The Recommended Architecture}

Level~1.5 is the compromise: constant learnable $\Pi$, per-token $R_t$.
This retains the per-token Kalman gain adaptation
\begin{equation}
  K_t = \frac{\Pi}{\Pi + R_t}
\end{equation}
but eliminates the M\"obius covariance scan. The state correction
\begin{equation}
  d_t = (1 - K_t) \, d_{t-1} + K_t \cdot \nu_t
\end{equation}
is a scalar affine recurrence with time-varying coefficient,
parallelizable via a single scalar Hillis-Steele scan: associative
combination
\begin{equation}
  (\alpha_1, v_1) \otimes (\alpha_2, v_2) = (\alpha_1 \alpha_2,\, \alpha_2 v_1 + v_2)
\end{equation}
yields cumulative affine transforms in $\mathcal{O}(\log T)$ depth,
$\mathcal{O}(T \log T)$ work (or $\mathcal{O}(T)$ with Mamba's
selective-scan kernel~\citep{gu2024mamba}).

Level~1.5 adds four learnable components to vanilla MapFormer:
\begin{enumerate}
  \item $\log \Pi \in \mathbb{R}^{H \times B}$: constant covariance
  \item $f_R$: MLP ($d_\mathrm{model} \to 128 \to H \cdot B$) producing
  per-token $\log R_t$
  \item $f_z$: MLP producing the measurement $\hat z_t$ (same as Level~1)
  \item No additional parameters in the scan itself
\end{enumerate}
Total: $\sim 50\mathrm{K}$ additional parameters, $\sim 10\%$ more than
vanilla MapFormer's 200K.

\subsection{Level 1.5-EM: Safe Initialization}

For the MapFormer-EM backbone (Hadamard product attention
$\mathrm{softmax}(A_X \odot A_P) \cdot V$), default initialization of
$\log R$ at zero gives $K{=}0.5$ at start. This corrupts $A_P$, which
Hadamard-products with $A_X$ to destroy gradient signal; 3 of 9 seeds
catastrophically diverge. We initialize $f_R$'s output bias at $3.0$,
giving $K \approx 0.05$ at start --- the InEKF behaves as a near-no-op
at init, then learns $R$ downward only where measurements are
informative. This is the residual-addition safe-init trick familiar from
ResNet's zero-init or adapter $\beta{=}0$.

\subsection{Predictive-Coding (PC) MapFormer}

An alternative correction mechanism based on forward models rather than
inverse. A forward model
\begin{equation}
  \hat o_t = g(\cos\theta_{\mathrm{path},t}, \sin\theta_{\mathrm{path},t})
\end{equation}
($2\pi$-invariant by construction) predicts the observation embedding
from position. The prediction error $\epsilon_t = x_t - \hat o_t$ is
mapped to a Lie-algebra correction $\delta\theta_t$ via an MLP, then
accumulated through a scalar affine scan with a learned gate. An
auxiliary loss $\|\epsilon\|^2$ at observation positions encourages the
forward model to actually model observations rather than collapse.

\subsection{Ablations}

Four Level~1.5 ablations isolate components:
\begin{description}
  \item[\texttt{L15\_ConstR}] constant learnable $R$ (no
  heteroscedasticity). Tests whether per-token $R_t$ matters.
  \item[\texttt{L15\_NoMeas}] measurement head output fixed to 0, so
  $\hat z = 0$ always. Tests whether the measurement signal contributes.
  \item[\texttt{L15\_NoCorr}] correction scan output forced to 0
  ($d_t = 0$). Recovers vanilla MapFormer with idle parameters. Sanity
  check.
  \item[\texttt{L15\_DARE}] $\Pi$ computed via scalar DARE from
  learnable $Q$ and a separate learnable $R_\mathrm{avg}$ (not
  learnable $\Pi$ directly). Tests whether the free learnable $\Pi$ adds
  expressivity beyond the DARE-prescribed value.
\end{description}
